%% ===========================================================================
\FloatBarrier
\section{Finite element formulation}
\label{sec:fem}
%% ===========================================================================

The weak formulation is discretized using a $C^1$ element, electrical
condensation and reduced-order fractional time integration. The element
matrices, solution procedure and numerical conditioning are described below.

\FloatBarrier
\subsection{Weak form}
\label{sec:weak}

Testing \eqref{eq:gov1} and \eqref{eq:gov2} against a virtual displacement
$\delta\bm{u}$ and a virtual potential $\delta\phi$ and integrating by parts
twice, the strain-gradient term produces both the double traction
\eqref{eq:doubletraction} and an integral over the element interiors,
\begin{multline}
  \int_\Omega \delta\varepsilon_{ij}\,C_{ijkl}\,\varepsilon_{kl}\,\mathrm{d}A
  + \tfrac12\int_\Omega \delta\varepsilon_{ij,m}\,M_{mn}\,C_{ijkl}\,
     \varepsilon_{kl,n}\,\mathrm{d}A
  - \int_\Omega \delta\varepsilon_{ij}\,e_{kij}\,E_k\,\mathrm{d}A \\
  - \int_\Omega \delta\varepsilon_{ij,m}\,\mu_{kijm}\,E_k\,\mathrm{d}A
  - \int_\Omega \delta E_k\left(e_{kij}\varepsilon_{ij}
     + \kappa_{kl}E_l + \mu_{kijl}\eta_{ijl}\right)\mathrm{d}A
  + \int_\Omega \rho\,\delta u_i\,\ddot u_i\,\mathrm{d}A
  = \delta W_{\rm body}
  + \int_{\Gamma}\left(\delta u_i\,t_i + \delta u_{i,j}n_j\,m_i\right)
    \mathrm{d}s ,
  \label{eq:weak}
\end{multline}
Here $W_{\rm body}$ denotes any prescribed body-force work; the surface
tractions and double tractions are displayed separately and must not be
counted again as part of $W_{\rm body}$. The reported vibration problems
are unforced in the interior. The displayed work assumes zero natural corner
forces (or constrained corner variations) and zero prescribed electrical
charge work. General data require the corresponding corner and electrical
boundary-work terms; floating-electrode zero net charge is imposed globally.
This is the basis of the assembled element
matrices. The boundary integral carries the classical and higher-order mechanical conditions of
Sec.~\ref{sec:bc_mech}; the electrical term enforces the prescribed electrode conditions. The
mechanical and electric blocks are assembled together, and electric
potential degrees of freedom are then eliminated by the electrode Schur
complement before the mechanical eigensolve or time integration.

\FloatBarrier
\subsection{The BFS element and its contributions}
\label{sec:bfs}

The domain is covered by Bogner--Fox--Schmit (BFS) rectangles. Each element has four
nodes and, per node, the eight mechanical degrees of freedom \eqref{eq:dofs}
carried by a bicubic Hermite tensor-product interpolation, that is 32
mechanical degrees of freedom per element, together with the four nodal
potentials of a bilinear interpolation. The strain and the strain gradient on
the element are
\begin{equation}
  \varepsilon = B(\xi,\eta)\,\bm{u}^e,\qquad
  \eta_k = \partial_k\varepsilon = B_{,k}(\xi,\eta)\,\bm{u}^e,
  \label{eq:Bmat}
\end{equation}
both integrated by $4\times4$ Gauss quadrature in the element frame, using the quadrature rule adopted for all elements; $3\times3$, $4\times4$ and
$5\times5$ rules change $k^2$ by at most $5.8\times10^{-6}$ and the fundamental
frequency by at most $7.1\times10^{-6}$ relative on the $48\times4$ mesh, supporting the selected quadrature order for this configuration. The gradient matrix $B_{,k}$ contains the second
derivatives of the shape functions and therefore requires $C^1$-conforming interpolation. The four element contributions are the elastic stiffness
\begin{equation}
  K_e=\int_{\Omega_e} B^{\mathsf T}C\,B\,\mathrm{d}A ,
  \label{eq:Ke}
\end{equation}
the gradient stiffness, in which the second moment replaces the scalar length,
\begin{equation}
  K_{g,e}=\tfrac12\int_{\Omega_e} M_{mn}B_{,m}^{\mathsf T}C\,B_{,n}\,\mathrm{d}A ,
  \label{eq:Kg}
\end{equation}
the electromechanical coupling, containing the piezoelectric and
single-constant flexoelectric contributions,
\begin{equation}
  B_e=\int_{\Omega_e}\left(B^{\mathsf T}e^{\mathsf T}B_\phi
       +\mu G_{\rm tr}^{\mathsf T}B_\phi\right)\,\mathrm{d}A ,
  \label{eq:Be}
\end{equation}
where $B_\phi=\nabla N_\phi$ (a $2\times4$ matrix) and the
$2\times32$ operator $G_{\rm tr}$ is defined by
$G_{\rm tr}\bm u^e=\nabla(\varepsilon_{xx}+\varepsilon_{yy})$.
This explicitly gives the tensor contraction implemented for the
isotropic one-constant flexoelectric specialization. The dielectric
and consistent mass matrices are
\begin{equation}
  C_{\phi,e}=\int_{\Omega_e}B_\phi^{\mathsf T}\kappa B_\phi\,\mathrm{d}A ,
  \label{eq:Cphi}
\end{equation}
and $M_e=\int_{\Omega_e}\rho N_u^{\mathsf T}N_u\,\mathrm{d}A$,
with $N_u$ the two-component mechanical interpolation matrix. In SI
units the material tensors have $[C]=\si{\pascal}$,
$[e]=\si{\coulomb\per\square\meter}$,
$[\kappa]=\si{\farad\per\meter}$,
$[\mu]=\si{\coulomb\per\meter}$ and $[M_{mn}]=\si{\square\meter}$.
The two-dimensional element integration gives energy per unit out-of-plane
width. If the generalized displacement components have units $[u_a]$, then
$[(K_e)_{ab}]=\si{\joule\per\meter}/([u_a][u_b])$,
$[(M_e)_{ab}]=\si{\kilogram\meter}/([u_a][u_b])$,
$[(B_e)_{aJ}]=\si{\coulomb\per\meter}/[u_a]$ and
$[(C_{\phi,e})_{IJ}]=\si{\farad\per\meter}$. These entries must not
be assigned a common unit when Hermite derivative DOFs are present.
Because the flexoelectric contribution involves $B_{,k}$, it is the term for which the
$C^1$ character of the element is indispensable.

\FloatBarrier
\subsection{Condensation of the potential}
\label{sec:condensation}

The assembled system is indefinite in the saddle-point form
\begin{equation}
  \begin{bmatrix} K & B \\ B^{\mathsf T} & -C_\phi \end{bmatrix}
  \begin{bmatrix} \bm{u} \\ \bm{\phi}\end{bmatrix}
  = \begin{bmatrix} \bm{f} \\ \bm{q}\end{bmatrix},
  \label{eq:saddle}
\end{equation}
where $\bm{u}$ and $\bm{\phi}$ collect nodal displacements and potentials,
$\bm{f}$ and $\bm{q}$ are mechanical and electrical loads, and $C_\phi$
is positive semidefinite before a reference potential is fixed (a constant
potential has zero electric field), and becomes positive definite on the
free electric unknowns once a grounded electrode removes that nullspace.
Eliminating the internal potential
degrees of freedom of the electrode set leaves the Schur complement
\begin{equation}
  K_{\rm cond} = K + B\,C_\phi^{-1}B^{\mathsf T},
  \qquad
  \bm{f}_{\rm cond} = \bm{f} + B\,C_\phi^{-1}\bm{q} ,
  \label{eq:schur}
\end{equation}
and eliminating the floating electrode as well reproduces
\eqref{eq:rankone}. Partition the free electric potential into interior
nodes (subscript $i$) and nodes on the single floating electrode
(subscript $f$), with the other electrode grounded; let $\bm g$
be the all-ones vector on the floating nodes. The actual block
condensation performed in the implementation gives
\begin{equation}
  d = (B_f-B_i C_{ii}^{-1}C_{if})\bm g ,\qquad
  \beta = \bm g^{\mathsf T}
  (C_{ff}-C_{fi} C_{ii}^{-1}C_{if})\bm g ,
  \label{eq:dbeta}
\end{equation}
where $B_i,B_f$ and $C_{ab}$ are the corresponding electrode-partitioned
subblocks of $B$ and $C_\phi$. The short-circuit stiffness includes
$B_iC_{ii}^{-1}B_i^{\mathsf T}$ before the rank-one update. The modal and quasi-static calculations use the pair $(d,\beta)$. The electrode
equations behind that pair are
\begin{equation}
  Q = d^{\mathsf T}\bm{u} - \beta V,
  \label{eq:electrode}
\end{equation}
with $Q$ the charge on the electrode that is held at the potential $V$ above
ground: short circuit ($V=0$) gives $Q=d^{\mathsf T}\bm{u}$, open circuit
($Q=0$) gives the sensed voltage
\begin{equation}
  V_{\rm oc} = \frac{d^{\mathsf T}\bm{u}}{\beta},
  \label{eq:voc}
\end{equation}
and $\beta=-\partial Q/\partial V$ at fixed displacement under the
charge-sign convention of \eqref{eq:electrode}, so $\beta>0$ is the
condensed capacitance magnitude of the electrode pair. For the two-dimensional model (quantities per unit out-of-plane width),
$Q$ has units \si{\coulomb\per\meter}, $V$ is a voltage and $\beta$
is \si{\farad\per\meter}. The components of $d$ have units of charge
per unit width divided by their corresponding generalized displacement:
\si{\coulomb\per\square\meter} for displacement DOFs measured in
meters, \si{\coulomb\per\meter} for first-derivative DOFs, and
\si{\coulomb} for mixed-derivative DOFs. A single unit for all entries
of the Hermite-coupling vector would be incorrect. The electrode charge and voltage are evaluated from $(d,\beta)$
and the nodal displacement without pointwise integration of the electric
displacement.

\FloatBarrier
\subsection{Modal L1 time stepping}
\label{sec:stepping}

With the L1 weights of \eqref{eq:l1}, both reported open-circuit
ring-down variants solve the same elastic problem but use different memory
operators:
\begin{equation}
  M\bm{a}^{n+1} + K_{oc}\bm{u}^{n+1} +
  \frac{\tau^\alpha K_d}{\Gamma(2-\alpha)}
  \sum_{k=0}^{n} b_{k}\,\frac{\bm{u}^{n+1-k}-\bm{u}^{n-k}}{\Delta t^{\alpha}}
  = \bm{f}^{n+1},\qquad K_d\in\{K_{oc},K_{uu}\},
  \label{eq:step}
\end{equation}
where $K_{oc}$ is the fixed open-circuit elastic stiffness and
$K_{uu}$ is the mechanical stiffness before electrical stiffening.
Define $K_{\rm mem}=\tau^\alpha K_d$, which has stiffness units
multiplied by $\mathrm{s}^{\alpha}$, so that the fractional derivative
of order $\alpha$ yields generalized forces. The terms ``memory on
$K_{oc}$'' and ``mechanical-only memory on $K_{uu}$'' denote distinct
phenomenological laws; neither is a dielectric-relaxation model. Because
the Hermite unknowns include displacements, slopes and mixed derivatives,
their mass and stiffness entries have coordinate-dependent SI units.
For the reported transient simulations the nodal displacement is projected
onto mass-normalized \emph{open-circuit} elastic modes,
\begin{equation}
  \bm{u}(t) \simeq \Phi\,\bm{q}(t),\qquad
  \Phi^{\mathsf T}M\Phi = I,\qquad
  \Phi^{\mathsf T}K_{oc}\Phi = \mathrm{diag}(\omega_{oc,i}^2),
  \label{eq:modal}
\end{equation}
and the convolution is carried on $\bm{q}$ instead of on the nodal vector, so the reduced linear system has the dimension of the retained modal
space. The direct L1 history sum nevertheless grows with the number of
elapsed time steps, and reconstruction of the physical displacement field
retains a dependence on the full mechanical degrees of freedom. Inserting the Newmark expressions for
$\bm{u}^{n+1}$ and $\bm{a}^{n+1}$ ($\beta_N=\tfrac14$, $\gamma_N=\tfrac12$),
\begin{equation}
  \bm{u}^{n+1} = \bm{u}^{n} + \Delta t\,\bm{v}^{n}
     + \Delta t^{2}\left[\left(\tfrac12-\beta_N\right)\bm{a}^{n}
       + \beta_N\,\bm{a}^{n+1}\right],
  \qquad
  \bm{v}^{n+1} = \bm{v}^{n} + \Delta t\left[(1-\gamma_N)\bm{a}^{n}
       + \gamma_N\,\bm{a}^{n+1}\right],
  \label{eq:newmark}
\end{equation}
gives a linear system for the unknown acceleration at step $n+1$,
\begin{equation}
  \left[M + \beta_N\Delta t^{2}K_{oc} +
    \frac{\beta_N\,\Delta t^{2-\alpha}}{\Gamma(2-\alpha)}K_{\rm mem}\right]
  \bm{a}^{n+1} = \bm{r}^{n+1},
  \label{eq:stepA}
\end{equation}
which is factorized once by Cholesky, the elastic and the memory stiffness both
appearing on the left; $\bm{r}^{n+1}$ collects the known displacements,
velocities and the history sum $\sum_{k\ge1}b_{k}\Delta\bm{u}^{n+1-k}$ of
\eqref{eq:l1}. For smooth initial data and $0<\alpha<1$, the Caputo history tends to
zero at the initial instant, giving
$\bm a^0=M^{-1}(\bm f^0-K_{oc}\bm u^0)$.
At $\alpha=1$, the initial derivative is the prescribed velocity, so
$\bm a^0=M^{-1}(\bm f^0-K_{oc}\bm u^0-\tau K_d\bm v^0)$.
The reported zero-initial-velocity ring-downs are unaffected by this distinction.

The time-discrete equation has two useful limits. At $\alpha=1$ the weights are $b_0=1$ and
$b_j=0$ for every $j\ge1$: the history vanishes identically, the operator
$\tau K_d\,\mathrm{D}^{1}$ becomes the backward difference $\tau
K_d(\bm{u}^{n+1}-\bm{u}^{n})/\Delta t$, and \eqref{eq:step} tends to
$M\bm{a}+K_{oc}\bm{u}+\tau K_d\bm{v}=\bm{f}$ up to
the $O(\Delta t)$ error of that difference, so the scheme is first order at
$\alpha=1$ as well. The computed ring-down reproduces the exact $\omega\tau/2$ to that
accuracy ($0.150111$, $0.150057$, $0.150029$ against $0.150000$
at $N$, $2N$ and $4N$ steps).
At $\tau\to0$ the memory matrix vanishes and the elastic problem is recovered.
The factor $\Delta t^{2-\alpha}K_{\rm mem}$ has the same
entry-wise units as $M$, since $K_{\rm mem}=\tau^\alpha K_d$ and
$[K_d]=[M]\,\mathrm{s}^{-2}$ for the corresponding generalized coordinates.
Thus the effective Newmark matrix is dimensionally homogeneous.

The truncation of the modal projection was measured in the reported
ring-down cases: with memory on $K_{oc}$ the motion is a
single mode and eight modes reproduce a forty-mode projection to below
$1\times10^{-12}$ of the initial voltage, a difference at the numerical round-off level, whereas with the memory on $K_{uu}$ the projected memory
couples the modes and eight modes differ from forty by $3.0\times10^{-4}$
(twenty modes: $4.7\times10^{-5}$). Eight modes are therefore
converged for the $\alpha=1$ decay quoted below, and the $\alpha<1$ decays are
read with that difference in mind. These finite mode-count comparisons apply only to the stated initial mode, electrodes and memory matrix; other responses require renewed truncation checks against a larger or full discrete reference.

\FloatBarrier
\subsection{Solution procedure}
\label{sec:algorithm}

For each geometry, assembly, condensation and solution proceed as follows.
\begin{enumerate}
\item Select the material tensor for the prescribed poling axis and
  reduce it to plane stress with
  \eqref{eq:Cp}--\eqref{eq:ep}.
\item Evaluate the second-moment tensor $M$ of the window from
  \eqref{eq:Mshapes}, rotated by \eqref{eq:Mrot2} if the window is not aligned.
\item Mesh the strip with $n_x\times n_y$ BFS rectangles and assemble
  $K_e$, $K_{g,e}$, $B_e$, $C_{\phi,e}$ and $M_e$ by $4\times4$ quadrature,
  equations \eqref{eq:Ke}--\eqref{eq:Cphi}.
\item Mark the electrode nodes of the chosen pair and split the potential
  degrees of freedom into free, earthed and floating.
\item Condense the potential, first to the electrode pair by \eqref{eq:schur}
  and then to the single floating electrode by \eqref{eq:dbeta}; store
  $(d,\beta)$.
\item \emph{Frequency domain:} solve the short-circuit pencil
  $K_{sc}\bm{\phi}=\omega^{2}M\bm{\phi}$ for the required number of modes
  (sparse shift--invert above $2600$ reduced unknowns, inverse generalized
  LAPACK eigensolve $M\bm{\phi}=\lambda K_{sc}\bm{\phi}$ for the
  smallest frequencies below that size, where $\omega^2=1/\lambda$),
  applying the same stiffness-diagonal congruence to both matrices.
\item Evaluate the coupling factor: the mode-based \eqref{eq:k2mode} when a
  mode is needed, the mode-free \eqref{eq:kstatic} when thicknesses or
  geometries are compared.
\item Post-process the open-circuit field: solve the internal potential from
  the condensed system at the prescribed electrode voltage, then the voltage
  from \eqref{eq:voc}.
\item \emph{Time domain:} project on the retained modes \eqref{eq:modal},
  factorize \eqref{eq:stepA} once, and march the L1 convolution with the
  weights \eqref{eq:bdef}; record the peaks of the modal coordinate.
\item Extract the damping ratio from two successive peaks with
  \eqref{eq:zeta}; for $\alpha<1$ report the first-period decrement instead.
\end{enumerate}

\FloatBarrier
\subsection{Conditioning and cost}
\label{sec:implementation}

Two properties of the implementation govern both cost and accuracy. The
ratio of gradient to classical elastic stiffness contributions scales as
$O(M/h^2)$ for comparable displacement variations over a transverse
length $h$, whereas the matrix entries associated with Hermite slope and
mixed-derivative degrees of freedom carry different dimensional scalings.
The assembled pencil can therefore become badly scaled as the thickness falls;
the implementation uses a common stiffness-diagonal congruence for both
matrices. In the dense low-mode route, the largest eigenvalues of the inverse
pencil $(M_s,K_s)$ yield the smallest physical frequencies; the resulting
modes are restored to the original coordinates and mass-normalized. This
reduces, but does not eliminate, conditioning sensitivity. For the $h=\SI{1}{\micro\meter}$ strip, the conditioned
stiffness matrix has a spectral condition number of order $10^{11}$ on tested
coarse meshes; the effect on the lowest frequency is noted in the numerical limitations. The $48\times4$ mesh contains $245$ nodes and $1960$ mechanical
degrees of freedom and is used for the main calculations; refining it to $96\times4$
changes the coupling factor by $5.9\times10^{-4}$ relative to the $48\times4$
value of the same clamped strip, and the thinnest strips are evaluated on this mesh because numerical
conditioning is the principal limitation. The assembly uses the strain and gradient matrices
\eqref{eq:Bmat}, the gradient stiffness \eqref{eq:Kg}, the coupling
\eqref{eq:Be} and the Schur complement \eqref{eq:saddle}, and the charge
constraint \eqref{eq:electrode} closes the electrode of the reduced model.
The mesh and time-step comparisons in Sec.~\ref{sec:verification} do not eliminate
conditioning sensitivity in the thinnest geometry.

\FloatBarrier
